\scanpage{19}
\deleted{\uline{§ \unsure{4|5} \unsure{Décompositions régulières}}}

\uline{§ 5} \uline{Localisation}

Soit $F$ satisfaisant a) b) c) d) e), ou encore $F_0$ satisfaisant $\alpha\,\beta\,\gamma\,\delta$ (c'est kif kif). Soit $A\in F$, considérons
\[
F_A=\{B\in\mathfrak P(A)\mid B\in F\}.
\]
Evidemment $F_A$ satisfait a) b) c), satisfait-il aussi d) et e)? Pour d), \unsure{cela|ceci} provient du fait immédiat que $F_A$ est engendré par $(F_A)_0$ i.e. que $\unsure{\forall A|U_A}\in F_A$ est réunion finie de parties loc fermées \inserted{$\in F_A$} de $A$. \deleted{\unsure{E}} Enfin, pour e) on considère
\[
A\supset A_1\supset A_2\supset\cdots\supset A_n\supset\cdots
\]
$A_i\in \unsure{F_A|F_{\deleted{A}}}$, $A_i$ rel. fermé dans $A$, $A_{i+1}$ rare dans $A_i$.
\margin{On peut partir d'une partie \uline{quelconque} $A$ de $X$, et poser $F_A=\{B\cap A\mid B\in F\}$; alors $F_A$ satisfait a) b) c) d) e)}

Alors $\overline{A_i}$ rare dans $\unsure{\overline{A_{i-1}}|\overline{A_i}}$, \deleted{\unsure{donc nulle pour $i$ grand on a}} (car $A_{i\deleted{\unsure{+1}}}-A_{i+1}$ dense dans $A_i$ dense dans $\overline{A_i}$, \deleted{donc $A_{i+1}$} \unsure{or|$\cap$}
\[
A_i-A_{i+1}=A_i-\underbrace{\overline{A_{i+1}}\cap A_i}_{A_{i+1}}\subset\overline{A_i}-\overline{A_{i+1}}
\]
qui est donc dense dans $\overline{A_i}$) donc pour $i$ grand on a $\overline{A_i}=\varnothing$ donc $A_i=\varnothing$ cqfd.

Prouvons que si $F$ satisfait aussi f), il en est de \unsure{$=$|même} de $F_A$. C'est trivial si $A$ loc fermé. Il reste à prouver que si $A\in F$ quelconque, ses comp. conn. sont $\in F$ et en nb fini. Or $A=\bigcup Z_i$, les $Z_i$ loc fermés, et par f) QPS les $Z_i$ \uline{connexes}. Donc les composantes
\mspage{20}
